\chapter{Contents of the attached archive}
\label{app:archive}

The reproduction archive shipped with the thesis contains the trained
models, the summary cross-validation metrics, and the gridded precipitation
product for the worked-example year~2023. Raw inputs (ReKIS station files,
the Copernicus DEM, SoilGrids rasters) and per-fold intermediate files are
excluded; they can be regenerated from the public sources by running the
pipeline scripts of the repository (Appendix~\ref{app:repo}).

\medskip\noindent\textbf{Note on availability.} The compressed archive
exceeds the $4$~GB upload limit of the faculty submission system and is
therefore not attached to this thesis directly. It is available on request;
the author can be contacted at \texttt{ismaktam@fit.cvut.cz}. The same
artefacts can also be regenerated by running the pipeline scripts of the
public repository against the project S3 bucket (Appendix~\ref{app:repo}).
The directory tree below documents the layout the archive follows.

\dirtree{%
.1 source/.
.2 README.md\DTcomment{layout, load snippets, list of omissions}.
.2 kriging/\DTcomment{Ordinary Kriging baseline (Chapter~\ref{ch:kriging})}.
.3 global\_variogram.pkl\DTcomment{fitted spherical and exponential variograms}.
.3 winner.json\DTcomment{selected transform $\times$ variogram combination}.
.3 holdout\_station\_ids.json\DTcomment{492 holdout station identifiers shared by all three methods}.
.3 kfold/cv\_results.json\DTcomment{5-fold MAE/CRPS summary}.
.3 test/.
.4 holdout\_predictions.parquet\DTcomment{492-station holdout predictions, 1961--2023}.
.4 test\_metrics.json\DTcomment{Table~\ref{tab:final-comparison} row}.
.2 lgbm/\DTcomment{Gradient boosting (Chapter~\ref{ch:gbm})}.
.3 final/.
.4 stage1.joblib\DTcomment{wet/dry binary classifier}.
.4 models.joblib\DTcomment{11 quantile regressors of the amount stage}.
.4 features.parquet\DTcomment{training feature matrix}.
.3 hparam/\DTcomment{HPO best parameters, grid results, sweep plots}.
.3 kfold/cv\_results.json\DTcomment{5-fold MAE/CRPS summary}.
.3 fold\_assignment.parquet\DTcomment{shared 5-fold split, used by all three methods}.
.3 oof\_predictions.parquet\DTcomment{out-of-fold predictions on training stations}.
.3 uncertainty\_fold0.json\DTcomment{quantile-spread diagnostics}.
.2 bayesnf/\DTcomment{Bayesian Neural Field (Chapter~\ref{ch:bayesnf})}.
.3 runs/vi\_\_final\_\_WY\_h1\_10\_\_ffrk\_full/.
.4 model.pkl\DTcomment{32-replica variational ensemble}.
.4 metrics.json.
.4 losses.npy\DTcomment{ELBO trace}.
.4 preds.parquet\DTcomment{holdout predictions}.
.3 final/features.parquet\DTcomment{training feature matrix}.
.3 kfold\_summary\_vi.json\DTcomment{5-fold VI summary}.
.3 kfold\_summary\_map.json\DTcomment{5-fold MAP comparison}.
.2 dataset/\DTcomment{multi-resolution gridded product, worked-example month (March 2023)}.
.3 grid\_\{2,5,10\}km\_static.parquet\DTcomment{cell geometry, elevation, soil; at 1\,km the static features are embedded in each monthly file}.
.3 grid\_\{1,2,5,10\}km/2023-03.parquet\DTcomment{input covariate stack for the worked example}.
.3 predictions\_\{1,2,5,10\}km/2023-03.parquet\DTcomment{predicted gridded fields}.
}

Loading each artefact is documented in \texttt{source/README.md}; in short,
the kriging variogram and the Bayesian Neural Field ensemble are unpickled,
the gradient-boosting boosters are read with \texttt{joblib.load}, and the
gridded product is read directly with \texttt{pandas.read\_parquet}. The
Python environment required to load the boosters and the Bayesian Neural
Field is installed from the repository root with
\texttt{pip install -e ".[cpu]"}.

\chapter{Source-code repository}
\label{app:repo}

The full source code, the experimental notebooks, the AWS execution
scripts, and the \LaTeX{} sources of the present document are released as a
public Git repository~\cite{thesisRepo}. The top-level layout is:

\dirtree{%
.1 precip\_interpolation\_thesis/.
.2 src/thesis/\DTcomment{installable Python package}.
.3 config.py\DTcomment{single source of truth for hyperparameters}.
.3 data/\DTcomment{ReKIS, DEM, SoilGrids loaders}.
.3 datasets/\DTcomment{per-day station dataset, prediction grid}.
.3 transforms/\DTcomment{Projection $\to$ Indicator $\to$ Detrend $\to$ NormalScore}.
.3 models/kriging/\DTcomment{two-stage Ordinary Kriging, variogram fitter, LOO-CV}.
.3 evaluation/\DTcomment{MAE/CRPS metrics, scoring helpers}.
.3 scripts/\DTcomment{pipeline entry points; one CLI per experiment}.
.2 notebooks/\DTcomment{experimental and pedagogical notebooks}.
.3 01\_data/\DTcomment{exploratory data analysis, dataset construction}.
.3 02\_methodology/\DTcomment{methodological derivations and demonstrations}.
.3 03\_kriging/\DTcomment{variogram fitting, LOO-CV, sensitivity}.
.3 04\_lgbm/\DTcomment{HPO sweep, 5-fold CV, final model, multi-resolution demo}.
.3 05\_bayesnf/\DTcomment{baseline, geo and time ablations, k-fold, final run}.
.3 06\_comparison/\DTcomment{cross-method evaluation against the held-out set}.
.3 07\_dataset/\DTcomment{multi-resolution gridded product generation}.
.2 scripts/\DTcomment{cloud orchestration (Docker, AWS, Vast.ai)}.
.3 entrypoint.sh\DTcomment{Docker entry point used on the GPU workers}.
.3 fetch\_source.sh\DTcomment{mirrors the curated S3 subset into source/}.
.2 thesis/\DTcomment{\LaTeX{} sources of the present document}.
.2 pyproject.toml\DTcomment{package metadata, CPU and GPU extras}.
.2 Dockerfile\DTcomment{CUDA 12.1, Python 3.12, uv}.
}

The pipeline is fully reproducible from this repository. After exporting
the AWS credentials documented in the repository \texttt{README}, the
script sequence
\texttt{run\_variogram}~$\to$~\texttt{build\_monthly\_grids}~$\to$~\texttt{run\_cv}
regenerates the kriging baseline. The notebook directories \texttt{04\_lgbm/}
and \texttt{05\_bayesnf/} regenerate the two machine-learning models, and
\texttt{07\_dataset/} regenerates the gridded product distributed in the
attached archive.
